\subsection{Conjunction}

\newcommand{\norwedge}[2]{{\nor}\norneg{#1}\norneg{#2}} % NOR hacking
\begin{notation}
    \[P \wedge Q \coloneqq \neg P \nor \neg Q \qquad [\norwedge{P}{Q}]\]
\end{notation}

% P, Q -> P and Q
\begin{align*}
    \premise{P}\\
    \premise{Q}\\
    \assumption[1]{\neg P}\\
    \byrule[1]{\bot}{\ref{rule:negation-elimination}} \\
    \assumption[1]{\neg Q}\\
    \byrule[1]{\bot}{\ref{rule:negation-elimination}}\\
    \byrule{\neg P \nor \neg Q}{\ref{rule:nor-introduction}}\\
    \bydefinition{P \wedge Q}
\end{align*}
\begin{equation}
    \inferrule{P \\ Q}{P \wedge Q}
    \tag{{$\wedge$}I}
    \label{rule:conjunction-introduction}
\end{equation}

% P and Q -> P, Q
\begin{align*}
    \premise{P \wedge Q}\\
    \bydefinition{\neg P \nor \neg Q}\\
    \assumption[1]{\neg P}\\
    \byrule[1]{\bot}{\ref{rule:nor-elimination}}\\
    \byrule{P}{\ref{rule:reductio-ad-absurdum}}
\end{align*}
\begin{align*}
    \premise{P \wedge Q}\\
    \bydefinition{\neg P \nor \neg Q}\\
    \assumption[1]{\neg Q}\\
    \byrule[1]{\bot}{\ref{rule:nor-elimination}}\\
    \byrule{Q}{\ref{rule:reductio-ad-absurdum}}
\end{align*}
% Therefore, wizardry?
\begin{equation}
    \inferrule{P \wedge Q}{P} \quad \inferrule{P \wedge Q}{Q}
    \tag{{$\wedge$}E}
    \label{rule:conjunction-elimination}
\end{equation}

%\begin{align*}
%    \premise{P \wedge Q}\\
%    \byrule{Q}{\ref{rule:conjunction-elimination}}\\
%    \byrule{P}{\ref{rule:conjunction-elimination}}\\
%    \byrule{Q \wedge P}{\ref{rule:conjunction-introduction}}
%\end{align*}
%%\newcommand{\wedgeC}{{\wedge}C}
%\[\inferrule{P \wedge Q}{Q \wedge P} \quad (\wedgeC)\]

%\begin{align*}
%    \premise{(P \wedge Q) \wedge R}\\
%    \byrule{P \wedge Q}{\ref{rule:conjunction-elimination}}\\
%    \byrule{P}{\ref{rule:conjunction-elimination}}\\
%    \byrule{Q}{\ref{rule:conjunction-elimination}}\\
%    \byrule{R}{\ref{rule:conjunction-elimination}}\\
%    \byrule{Q \wedge R}{\ref{rule:conjunction-introduction}}\\
%    \byrule{P \wedge (Q \wedge R)}{\ref{rule:conjunction-introduction}}
%\end{align*}
%%\newcommand{\wedgeA}{{\wedge}A}
%\[\inferrule{(P \wedge Q) \wedge R}{P \wedge (Q \wedge R)} \quad (\wedgeA)\]
